\section{Axiom of Choice}




\begin{definition}
    \label{defChoiceFunction}
    %------------------------------------------------------------
    A function $f$ is a \emph{choice function} for a set $A$ if it maps every
    nonempty subset $S \subseteq A$ into an element of the set $f(S) \in S$.
\end{definition}




\begin{definition}
    \label{defUpperChoiceFunction}
    %------------------------------------------------------------
    Let $(P, \le)$ be a partial order, $A \subseteq P$ a subset of the order and
    $f$ a choice function for $P$. Then the function $g$
    $$
        g(A) := f(A^u - A)
    $$
    defined for all $A$ such that $A^u - A$ is not empty is an \emph{upper bound
    choice function} for the partial order. It selects for all subsets of the
    order a strict upper bound of the set, provided that a strict upper bound
    exists.
\end{definition}


\begin{definition}
    \label{defInitialSegment}
    %------------------------------------------------------------
    Let $(P, \le)$ be a partial order and $C \subseteq P$ a chain in $P$. $A$ is
    an \emph{initial segment of the chain $C$} if $A \subseteq C$ and
    $$
    \rulev{
        x \in A
            \\
            y \le x
            \\
            y \in C
    } {
        y \in A
    }
    $$
    is valid.
\end{definition}




In the remainder of the section we assume a partial order $(P,\le)$ with an
upper bound choice function $g$ which selects from every subset $A \subseteq P$
a strict upper bound i.e. an element of $A^u - A$ provided that $A^u - A$ is not
empty i.e. a strict upper bound exists.


\begin{definition}
    \label{defTower}
    %------------------------------------------------------------
    $T \subseteq P$ is a \emph{tower} if it is a chain and for all proper inital
    segments~\ref{defInitialSegment} $A \subset T$ the value $g(A)$ is the least
    strict upper bound of $A$ in $T$.
\end{definition}


\begin{lemma}
    \label{thmTowerWellorder}
    \emph{A tower is a wellorder.}

    \begin{proof}
        We have to show that each nonempty subset $S \subseteq T$ of a tower $T$
        has a least element. We form the strict lower bounds $A := (S^l - S)
        \cap T$ of $S$ in the tower $T$ and prove that $S$ has a least element
        in steps by proving
            $A$ is an initial segment of $T$,
            $A$ is a proper initial segment of $T$,
            $g(A) \in S$
            and
            $g(A)$ is the least element of $S$


        \begin{enumerate}
            \item $A$ is an initial segment of $T$: By definition $A$ is a
                subset of $T$. Let $x \in A$, $y \le x$ and $y \in T$. We have
                to show that $y \in A$ is valid. Since $A$ is downclosed in $T$
                this is certainly valid.

            \item $A$ is a proper initial segment of $T$: By definition of $A$
                the sets $A$ and $S$ are disjoint and their union is a subset of
                $T$.  Since $S$ is not empty, $A$ is proper initial segment of
                $T$.

            \item $g(A) \in S$:
            By definition of  a tower~\ref{defTower}
            $g(A)$ is the least strict upper bound of $A$ in $T$. If it were not in
            $S$, then $g(A)$ would be a strict lower bound of $S$ in $T$ and would
            have to be in $A$ by the definition of $A$. This contradicts the fact
            that $g(A)$ is a strict upper bound of $A$ in $T$.

            \item $g(A)$ is the least element of $S$:
            $S$ has only strict upper bounds of $A$ by definition of $A$, $g(A)$ is
            the least strict upper bound and it is in $S$. Therefore $g(A)$ is the
            least element of $S$.
        \end{enumerate}
    \end{proof}
\end{lemma}



\begin{lemma}
    \label{thmTowerInitialSegmentEnlargement}
    \emph{Let $T$ be a tower and $I \subset T$ a proper initial segment of $T$. Then
    $A \cup \set{g(I)}$ is an initial segment of $T$}.

    \begin{proof}
        Since $g(I)$ selects the least strict upper bound of $I$ in $T$ we have
        $I \cup \set{g(I)} \subseteq T$ and therefore satisfies the first
        condition in the definition of an initial
        segment~\ref{defInitialSegment}.

        For the second condition assume $x \in I \cup \set{g(I)}$, $y \le x$ and
        $y \in T$. We have to show that $y \in I \cup \set{g(I)}$.

        We distinguish the cases $x \ne g(I)$ and $x = g(I)$. In the first case
        $x \in I$ and therefore $y \in I$ (which implies $y \in I \cup
        \set{g(I)}$) because $I$ is an initial segment of $T$.

        Now we consider the case $x = g(I)$. For both cases $y = x$ and $y < x$
        we have $y \in I \cup \set{g(I)}$.
    \end{proof}
\end{lemma}



\begin{lemma}
    \label{thmTowersComparable}
    %------------------------------------------------------------
    \emph{
        Comparability Lemma:
        Let $B$ and $C$ be towers~\ref{defTower}. Then either $B$ is an initial
        segment of $C$ or $C$ is an initial segment of $B$.
    }

    \begin{proof}
        Let $A := \set{x \in B \cap C: \set{x}^l \cap B = \set{x}^l \cap c}$.

        \begin{enumerate}
            \item $A$ is an initial segment of $B$:

                We assume that $x \in A$ $y \in B$ and $y \le x$ and have to
                show $y \in A$.

                Since $x \in A$ then by definition $\set{x}^l \cap B = \set{x}^l
                \cap C$.

                We have $y \le x$ and $y \in B$. Therefore $y \in \set{x}^l \cap
                B$. This implies $y \in \set{x}^l \cap C$ and therefore $y \in
                C$.

                In order to show $y \in A$ we have to prove $\set{y}^l \cap B =
                \set{x}^l \cap C$. We first show $\set{y}^l \cap B \subseteq
                \set{x}^l \cap C$.

                Assume $z \in \set{x}^l \cap B$. Because of $y \le x$ we have $z
                \in \set{x}^l \cap B$ and therefore $z \in \set{x}^l \cap C$
                i.e. $z \in C$. And since $z \in \set{y}^l$ we have $z \in
                \set{y}^l \cap C$.

                The other direction can be proved in the same manner.

            \item $A$ is an initial segment of $C$: Proof similar to the previous
                step.

            \item Every initial segment $I$ of $B$ and $C$ is a subset of $A$:

                Assume $x \in I$. We have to show $x \in B \cap C$ and
                $\set{x}^l \cap B = \set{x}^l \cap C$.

                The first statement is trivial since $I$ is a common initial
                segment of $B$ and $C$.

                In order to prove the second statement we show $\set{x}^l \cap B
                \subseteq \set{x}^l \cap C$. The other direction is similar.

                Assume $y \in \set{x}^l \cap B$ i.e. $y \le x$ and $y \in B$.
                This implies $y \in I$ because $I$ is an initial segment of $B$
                by definition of an initial segment~\ref{defInitialSegment}.
                Since $I$ is an initial segment of $C$ as well by definition of
                an initial segment we have $y \in C$. Therefore $y \in \set{x}^l
                \cap C$.

            \item $A$ is not a proper initial segment of both $B$ and $C$:

                Assume that $A$ is a proper initial segment of $B$ and $C$.
                Because $B$ and $C$ are towers~\ref{defTower}
                and~\ref{thmTowerInitialSegmentEnlargement} the set $A \cup
                \set{g(A)}$ is an initial segment of $B$ and $C$. This
                contradicts the fact the every initial segment of $B$ and $C$ is
                a subset of $A$.

            \item Either $B$ is an initial segment of $C$ or $C$ is an initial
                segment of $B$:

                Since $A$ is not a proper initial segment of $B$ and $C$ but an
                initial segment of both it has to be identical to one of them.
                So one of $B$ and $C$ is an initial segment of the other.
        \end{enumerate}
    \end{proof}
\end{lemma}



\begin{lemma}
    \label{thmTowerUnion}
    %------------------------------------------------------------
    \emph{
        The union $U$ of all towers is a tower.
    }

    \begin{proof}
        We have to prove that $U$ is a chain and for all proper initial segments
        $A \subset U$ of $U$ the upper bound choice function applied to the
        initial segment $g(I)$ selects the least strict upper bound of $I$ in
        $U$.

        \begin{enumerate}
            \item $U$ is a chain: Let $x,y \in U$. We have to show that either
            $x \le y$ or $y \le x$ is valid.

            Since $x,y \in U$ there are two towers $B$ and $C$ such that $x \in
            B$ and $y \in C$ by the definition of a union. By the comparability
            lemma~\ref{thmTowersComparable} we know that one of $B$ and $C$ is a
            ideal of the other. Without loss of generality we assume $B$ is an
            initial segment of $C$. This implies that $x \in C$ is valid.

            Since $C$ is a tower and therefore a chain either $x \le y$ or $y
            \le x$ is valid.


            \item Strict lower bound selected:

                Assume $I$ is a proper initial segment of $T$. We have to show
                that $g(I)$ selects the least strict upper bound of $I$ in $T$.

                MISSING!!
        \end{enumerate}
    \end{proof}
\end{lemma}
